\documentclass{article}

\usepackage{geometry}
\geometry{a4paper, margin=2.5cm}

\usepackage{amsmath, amssymb, amsthm}
\usepackage{algorithm}
\usepackage{algpseudocode}
\usepackage{booktabs}
\usepackage{xcolor}

\theoremstyle{definition}
\newtheorem{lemma}{Lemma}
\newtheorem{theorem}{Theorem}

\title{Design and Analysis of Directed Minimum Spanning Tree Algorithm}
\author{}
\date{\today}

\begin{document}

\maketitle

% ============================================================
% Section 1: Algorithm Description
% ============================================================
\section{Algorithm Description}

\subsection{Algorithm Overview}

This algorithm solves the Directed Minimum Spanning Tree (DMST) problem, also known as the minimum arborescence problem. We build upon the Chu-Liu/Edmonds algorithm and optimize it using a \textbf{mergeable heap} (leftist tree), \textbf{disjoint set union} (DSU), and \textbf{lazy propagation} to achieve $O(n \log m)$ time complexity.

\subsection{Data Structures}

\subsubsection{Leftist Heap (Mergeable Heap)}

The leftist heap is a variant of the binary heap that supports efficient merging. For each node $v$ in the graph, we maintain a leftist heap $H[v]$ storing all incoming edges to $v$.

\paragraph{Heap Node Structure.} Each heap node stores:
\begin{itemize}
    \item $id$: identifier of the original edge in the graph
    \item $w$: edge weight (stored as original value, actual weight is $w + add$)
    \item $dist$: the distance of the node, defined as the length of the shortest path to a null child
    \item $l, r$: pointers to left and right children
\end{itemize}

\paragraph{Leftist Property.} For every node in the heap, the distance of the left child is at least the distance of the right child ($dist(left) \geq dist(right)$). This ensures the right spine has length $O(\log n)$.

\paragraph{Merge Operation.} The merge of two leftist heaps $a$ and $b$ is performed recursively:
\begin{enumerate}
    \item If either heap is empty, return the other.
    \item Ensure $a.w \leq b.w$ by swapping if necessary (min-heap property).
    \item Recursively merge $a.r$ (right child) with $b$, and assign the result to $a.r$.
    \item If $dist(a.l) < dist(a.r)$, swap $a.l$ and $a.r$ (maintain leftist property).
    \item Update $dist(a) = dist(a.r) + 1$ (or 0 if right child is null).
    \item Return $a$.
\end{enumerate}
The merge operation runs in $O(\log m)$ time due to the leftist property bounding the right spine length.

\paragraph{Pop Operation.} To remove the minimum element, we simply merge the left and right children: $pop(h) = merge(h.l, h.r)$. This also runs in $O(\log m)$ time.

\subsubsection{Other Data Structures}

\begin{itemize}
    \item \textbf{Disjoint Set Union} $dsu$: Tracks contracted super-nodes. Each set represents a contracted cycle, with the representative being the "representative element."
    \item \textbf{Lazy Tag} $add[v]$: Stores the offset for edge weights in $H[v]$. The actual edge weight is $w(e) + add[v]$, but we store original weights in the heap to avoid frequent updates.
    \item \textbf{Predecessor Array} $pre[v]$: Records the selected incoming edge for node $v$.
\end{itemize}

\subsection{Algorithm Workflow}

\begin{algorithm}[H]
\caption{Chu-Liu/Edmonds Algorithm with Mergeable Heap}
\label{alg:dmst}
\begin{algorithmic}[1]
\Require Directed graph $G=(V, E)$, root node $root$, edge weights $w: E \to \mathbb{R}$
\Ensure Weight of minimum arborescence

\State Initialize leftist heap $H[v]$ for each $v \in V$ with all incoming edges
\State Initialize $dsu$, $pre[v] \gets -1$, $vis[v] \gets -1$, $add[v] \gets 0$
\State $answer \gets 0$

\For{$i = 0$ \textbf{to} $|V|-1$}
    \If{$i = root$}
        \State \textbf{continue}
    \EndIf
    
    \State $u \gets i$, $path \gets \emptyset$
    
    % Traverse along predecessor chain
    \While{$vis[u] \neq i$ \textbf{and} $u \neq root$ \textbf{and} $dsu.find(u) = u$}
        \State $vis[u] \gets i$, $path \gets path \cup \{u\}$
        
        % Get minimum valid incoming edge
        \While{$H[u] \neq \emptyset$}
            \State $e \gets H[u].top()$
            \If{$dsu.find(e.from) = dsu.find(u)$}
                \State $H[u].pop()$ \Comment{Skip edges within same component}
            \Else
                \State \textbf{break}
            \EndIf
        \EndWhile
        
        \If{$H[u] = \emptyset$}
            \State \Return $-1$ \Comment{No valid incoming edge}
        \EndIf
        
        \State $add[u] \gets w(e) - H[u].min()$ \Comment{Update lazy tag}
        \State $pre[u] \gets e.from$
        \State $answer \gets answer + w(e) + add[u]$
        \State $u \gets pre[u]$
    \EndWhile
    
    \If{$u = root$ \textbf{or} $dsu.find(u) \neq u$}
        \State \textbf{continue} \Comment{Connected to existing structure}
    \EndIf
    
    % Contract cycle
    \State $cur \gets dsu.find(pre[u])$
    \For{$v \in path$}
        \If{$dsu.find(v) = cur$}
            \State \textbf{break}
        \EndIf
        \State $dsu.f[v] \gets cur$
        \State Push down lazy tags and merge $H[cur]$ with $H[v]$
    \EndFor
    
    \While{$H[cur] \neq \emptyset$ \textbf{and} $dsu.find(H[cur].top().to) = cur$}
        \State $H[cur].pop()$ \Comment{Remove edges pointing into cycle}
    \EndWhile
    
    \State $i \gets i-1$ \Comment{Re-process the representative}
\EndFor

\State \Return $answer$
\end{algorithmic}
\end{algorithm}

\subsection{Key Design Decisions}

\paragraph{Why stop when encountering a non-representative element?} 
When traversing the predecessor chain, if we reach a node $u$ where $dsu.find(u) \neq u$, it means $u$ has been merged into a super-node. This indicates that $u$ is part of a previously processed structure that is guaranteed to eventually connect to the root (either directly or through further contractions). Therefore, we can safely stop traversing.

\paragraph{Lazy propagation mechanism.}
We maintain $add[v]$ as the accumulated offset for heap $H[v]$. When accessing the minimum edge, the actual weight is $H[v].min() + add[v]$. When merging two heaps, we push down the lazy tags first. This avoids $O(m)$ explicit edge weight updates during contractions.

\paragraph{Cycle contraction.}
When a cycle is detected (we encounter a node already visited in the current iteration), we contract all nodes in the cycle into a single super-node. The heaps are merged, and edges pointing into the cycle are lazily removed. The algorithm then re-processes this super-node to find its incoming edge from outside the cycle.

% ============================================================
% Section 2: Complexity Analysis
% ============================================================
\section{Complexity Analysis}

\subsection{Heap Initialization}

We insert each of the $m$ edges into the heap of its destination node. For a leftist heap, insertion is equivalent to a merge operation with a single-node heap, which takes $O(\log m)$ time per edge. Therefore, the total initialization time is:
\[
T_{\text{init}} = O(m \log m)
\]

\subsection{Main Loop Analysis}

The outer loop iterates over all $n$ nodes. For each node, we traverse its predecessor chain until one of three conditions is met:
\begin{enumerate}
    \item We reach the root;
    \item We encounter a node already visited in the current iteration (cycle detected);
    \item We encounter a node that is not a representative ($dsu.find(u) \neq u$).
\end{enumerate}

\paragraph{DSU operations.} Each $find$ or $union$ operation takes $O(\alpha(n))$ amortized time, where $\alpha$ is the inverse Ackermann function, effectively constant.

\paragraph{Heap operations during traversal.} Each edge is popped from a heap at most once (when its source and destination are in the same component). Therefore, the total number of $pop$ operations across the entire algorithm is $O(m)$.

\paragraph{Cycle contraction.} Each contraction reduces the number of DSU components by at least one. Since we start with $n$ components and end with at least one, there are at most $n-1$ contractions. Each contraction involves merging heaps of the cycle nodes. The total cost of all heap merges is $O(n \log m)$ because each node participates in at most $O(\log n)$ merges (similar to union-find analysis), and each merge costs $O(\log m)$.

\paragraph{Re-processing nodes.} When we contract a cycle, we decrement $i$ to re-process the representative. However, each node becomes a non-representative (gets merged) at most once. Therefore, the condition $dsu.find(u) \neq u$ prevents further processing, and the total number of iterations where actual work is done is $O(n)$.

\subsection{Overall Complexity}

Combining all components:
\begin{align*}
T(n, m) &= T_{\text{init}} + T_{\text{traversal}} + T_{\text{contraction}} \\
&= O(m \log m) + O(n \cdot \alpha(n)) + O(n \log m) \\
&= O(m \log m) + O(n \log m)
\end{align*}

In a connected graph, $O(m \log m) = O(m \log n)$. Therefore, the overall time complexity is:
\[
\boxed{O(m \log n)}
\]

We then analysis the space complexity. We maintain a leftist heap for each node, which can store up to $O(m)$ edges in total across all heaps. The DSU structure and other auxiliary arrays take $O(n)$ space. Therefore, the overall space complexity is:
\[O(n + m)\]


\end{document}
